\section{Exercises}\label{sec:08-group-theorems:05-exercises:1}

\textbf{Key points.} Three facts now hold for \emph{every} group, the identity is unique, a left inverse equals the (unique) two-sided inverse, and \((ab)^{-1}=b^{-1}a^{-1}\). The recurring proof pattern is relating both sides of a goal to a common third expression, or padding with the identity and cancelling. Once a uniqueness fact is proved, later goals can be \emph{characterized} by it instead of computed directly.

\begin{enumerate}
\def\labelenumi{\arabic{enumi}.}
\tightlist
\item
  The proof of \texttt{id\_\allowbreak{}unique} related \texttt{e'} and \texttt{Grp.id} to a common third expression, \texttt{Grp.op e' Grp.id}. Explain whether the same proof could have gone through by relating them to \texttt{Grp.op Grp.id e'} instead, given that \texttt{h} only says something about \texttt{Grp.op e' a} for arbitrary \texttt{a}, and nothing about \texttt{Grp.op a e'}.
\item
  \texttt{inv\_\allowbreak{}op} proved \((ab)^{-1}=b^{-1}a^{-1}\) by showing \(b^{-1}a^{-1}\) is a left inverse of \(ab\), rather than computing \((ab)^{-1}\) directly. Explain whether the axioms of \texttt{Group} offer any way to compute it directly, and if not, why not.
\item
  Prove \texttt{theorem inv\_\allowbreak{}inv (a : G) : Grp.inv (Grp.inv a) = a}. Before writing any tactics, consider whether this matches the shape of a lemma already in hand (Theorem 2 again). What single fact about \texttt{a} and \texttt{Grp.inv a} would permit invoking it directly?
\item
  Prove \texttt{theorem cancel\_\allowbreak{}left (a b c : G) (h : Grp.op a b = Grp.op a c) : b = c}. \texttt{b} and \texttt{c} cannot be rewritten directly in isolation, they only appear inside \texttt{h}. This uses the same ``regroup, then cancel'' pattern as Theorem 3.
\end{enumerate}

Solutions, \hyperref[sec:15-appendix-solutions:08-chapter-8:1]{Appendix, Chapter 8}.

